%! AP Statistics — Unit 2, Lesson 2.6 — Group Activity
\documentclass[10pt]{article}
\usepackage{apstats-article}
\usepackage{apstats-boxes}

\begin{document}

\pageheader{Unit 2, Lesson 2.6}{Group Activity}
\namepartnerperiod

\vspace{0.18in}

% ── SCENARIO ─────────────────────────────────────────────────────────────────
\begin{scenariobox}[Context]{sky}
\small
Each tier uses a different regression equation from a real-world context. Your group will
interpret the equation, make predictions, and assess any limitations. Write all answers in
complete sentences with context.
\end{scenariobox}

\vspace{0.12in}

% ── TIER R ────────────────────────────────────────────────────────────────────
\begin{tcolorbox}[colback=white, colframe=black!40, boxrule=0.8pt, arc=2mm,
    left=4mm, right=4mm, top=3mm, bottom=3mm,
    title={\textbf{Tier R --- Remediate}}, fonttitle=\bfseries, halign title=center]
\small

\textbf{Context: Coffee Shop Sales}

A coffee shop owner collected data on daily high temperature (°F) and the number of hot
coffee drinks sold that day. Technology produced the following regression equation:
\[
  \widehat{\text{coffee}} = 210 - 2.4(\text{temp})
\]
The data were collected for days with temperatures ranging from \textbf{30°F to 75°F}.

\vspace{10pt}

\begin{enumerate}[leftmargin=1.6em, itemsep=0.2in]

  \item Identify the explanatory and response variables.\\[6pt]
        Explanatory: \blank{5.5cm} \quad Response: \blank{5.5cm}

  \item Interpret the \textbf{slope} ($-2.4$) in context.\\[8pt]
        \writeline\\[3pt]\writeline\\[3pt]\writeline

  \item Interpret the \textbf{$y$-intercept} ($210$) in context.
        Is it meaningful? Explain.\\[8pt]
        \writeline\\[3pt]\writeline\\[3pt]\writeline

  \item Predict the number of coffee drinks sold on a day when the temperature is \textbf{55°F}.
        Show your work.\\[8pt]
        \writeline\\[3pt]\writeline

\end{enumerate}
\end{tcolorbox}

\vspace{0.12in}

% ── TIER A ────────────────────────────────────────────────────────────────────
\begin{tcolorbox}[colback=white, colframe=black!40, boxrule=0.8pt, arc=2mm,
    left=4mm, right=4mm, top=3mm, bottom=3mm,
    title={\textbf{Tier A --- Approaching Proficiency}}, fonttitle=\bfseries, halign title=center]
\small

\textbf{Context: Sleep and Reaction Time}

A researcher recorded hours of sleep ($x$) and reaction time in milliseconds ($y$) for
20 adults. The data ranged from \textbf{4 to 9 hours} of sleep. Technology output:
\[
  \widehat{\text{RT}} = 430 - 25(\text{sleep})
\]

\vspace{8pt}

\begin{enumerate}[leftmargin=1.6em, itemsep=0.18in]

  \item Interpret the \textbf{slope} ($-25$) in context.\\[8pt]
        \writeline\\[3pt]\writeline\\[3pt]\writeline

  \item Interpret the \textbf{$y$-intercept} ($430$) in context.
        Is it meaningful? \textbf{Support your answer.}\\[8pt]
        \writeline\\[3pt]\writeline\\[3pt]\writeline

  \item Predict the reaction time for someone who slept \textbf{7 hours}.
        Is this interpolation or extrapolation? Show work.\\[8pt]
        \writeline\\[3pt]\writeline

  \item Predict the reaction time for someone who slept \textbf{1 hour}.
        Is this reliable? Why or why not?\\[8pt]
        \writeline\\[3pt]\writeline

\end{enumerate}
\end{tcolorbox}

\vspace{0.12in}

% ── TIER E ────────────────────────────────────────────────────────────────────
\begin{tcolorbox}[colback=white, colframe=black!40, boxrule=0.8pt, arc=2mm,
    left=4mm, right=4mm, top=3mm, bottom=3mm,
    title={\textbf{Tier E --- Extension}}, fonttitle=\bfseries, halign title=center]
\small

\textbf{Context: Wrist Circumference and Height}

A medical researcher recorded wrist circumference in cm ($x$) and height in cm ($y$) for
50 adults. Wrist measurements ranged from \textbf{14 cm to 20 cm}. Output:
\[
  \widehat{\text{height}} = -1.2 + 10.4(\text{wrist})
\]

\vspace{8pt}

\begin{enumerate}[leftmargin=1.6em, itemsep=0.18in]

  \item Interpret the \textbf{slope} in context.\\[8pt]
        \writeline\\[3pt]\writeline\\[3pt]\writeline

  \item Interpret the \textbf{$y$-intercept} in context.
        Is it meaningful? Carefully support your answer --- this one is tricky.\\[8pt]
        \writeline\\[3pt]\writeline\\[3pt]\writeline

  \item Predict the height of a person with a 17 cm wrist circumference.
        Classify as interpolation or extrapolation and assess reliability.\\[8pt]
        \writeline\\[3pt]\writeline

  \item A colleague suggests using this equation to predict the height of a person with a
        0 cm wrist circumference. What prediction results? What does this tell you about
        the limitation of the $y$-intercept interpretation?\\[8pt]
        \writeline\\[3pt]\writeline\\[3pt]\writeline

\end{enumerate}
\end{tcolorbox}

\end{document}
